\chapter{Testing, Experimentation, and Performance Evaluation}
\label{ch:testing}

This chapter describes how the implemented system is tested and evaluated. Each subsystem is tested independently first, and then the components are tested together in an integrated end-to-end scenario. The evaluation covers functional correctness, security properties, AI moderation accuracy, IPFS storage performance, queue reliability, and blockchain network performance.

\section{Experimental Setup and Evaluation Environment}
\label{sec:exp-setup}

\subsection{Hardware Infrastructure}

All tests are performed on cloud-based VPS infrastructure. Two servers are used to match the production deployment configuration:

\begin{table}[h!]
\centering
\caption{VPS Hardware Specifications for the Evaluation Environment}
\label{tab:hardware-specs}
\begin{tabular}{|l|l|}
\hline
\textbf{Specification} & \textbf{Value} \\
\hline
CPU & 4 vCPU cores \\
\hline
RAM & 8 GB \\
\hline
Storage & 75 GB SSD \\
\hline
Network & High-bandwidth cloud uplink \\
\hline
Operating System & Ubuntu Server 22.04 LTS \\
\hline
\end{tabular}
\end{table}

Server 1 hosts the core operational services: the private PoA blockchain, the backend relayer (with Redis), the ZKP GovID Simulator, and the web DApp. Server 2 hosts the AI oracle ensemble (four Docker containers) and the IPFS storage service.

\subsection{Software and Deployment Configuration}

All services are containerized using Docker. Dokploy is used as the self-hosted deployment platform on both servers. CI/CD pipelines are configured via GitHub webhooks to the \texttt{develop} branch. Smart contracts are compiled with Hardhat using Solidity 0.8.20. Automated contract testing uses the Hardhat test runner with the Chai assertion library.

\section{Subsystem Unit and Integration Testing}
\label{sec:unit-testing}

\subsection{Smart Contract Testing (Reporting.sol)}

Smart contract tests are run using the Hardhat framework with Chai assertions on a local in-process blockchain. Table~\ref{tab:reporting-tests} lists all test cases for \texttt{Reporting.sol}, organized by functional group.

\begin{table}[h!]
\centering
\caption{Reporting.sol Hardhat Test Suite (T1.1.1 -- T1.1.41)}
\label{tab:reporting-tests}
\begin{tabular}{|p{1.3cm}|p{10cm}|p{1.5cm}|}
\hline
\textbf{Test ID} & \textbf{Description} & \textbf{Result} \\
\hline
\multicolumn{3}{|l|}{\textit{Group 1 — Report Submission (submitReport)}} \\
\hline
T1.1.1 & Authorized relayer submits a report; \texttt{ReportSubmitted} event is emitted with correct fields & Passed \\
\hline
T1.1.2 & Non-relayer address is rejected with \texttt{Unauthorized()} & Passed \\
\hline
T1.1.3 & Duplicate nullifier is rejected with \texttt{NullifierAlreadyUsed()} & Passed \\
\hline
T1.1.4 & Empty CID is rejected with \texttt{EmptyCID()} & Passed \\
\hline
T1.1.5 & Zero \texttt{reportHash} is rejected with \texttt{InvalidHash()} & Passed \\
\hline
T1.1.6 & Zero nullifier is rejected with \texttt{InvalidNullifier()} & Passed \\
\hline
T1.1.7 & Zero pseudonym is rejected with \texttt{InvalidPseudonym()} & Passed \\
\hline
T1.1.8 & Report count increments correctly across multiple submissions & Passed \\
\hline
T1.1.9 & \texttt{phaseDeadline} is set on submission & Passed \\
\hline
\multicolumn{3}{|l|}{\textit{Group 2 — Validation Voting (castValidationVote)}} \\
\hline
T1.1.10 & Upvote increments \texttt{validationUpvotes} and emits \texttt{ValidationVoteCast} & Passed \\
\hline
T1.1.11 & Downvote increments \texttt{validationDownvotes} & Passed \\
\hline
T1.1.12 & Same vote nullifier is rejected with \texttt{NullifierAlreadyUsed()} & Passed \\
\hline
T1.1.13 & Same citizen pseudonym cannot vote twice (\texttt{CitizenAlreadyVoted}) & Passed \\
\hline
T1.1.14 & Vote after window expiry triggers lazy finalization (late vote not counted) & Passed \\
\hline
T1.1.15 & Voting reverts when report is not in \texttt{PendingValidation} state & Passed \\
\hline
\multicolumn{3}{|l|}{\textit{Group 3 — Voting Window Finalization}} \\
\hline
T1.1.16 & Upvotes $>$ downvotes after window $\Rightarrow$ status becomes \texttt{Open} (2) & Passed \\
\hline
T1.1.17 & Downvotes $\geq$ upvotes (including zero votes) $\Rightarrow$ \texttt{CommunityRejected} (1) & Passed \\
\hline
T1.1.18 & \texttt{finalizeVotingWindow} reverts when window is still open & Passed \\
\hline
T1.1.19 & \texttt{finalizeVotingWindow} succeeds exactly after window duration (time-travel test) & Passed \\
\hline
T1.1.20 & \texttt{batchFinalizeVotingWindows} finalizes multiple expired reports in one call & Passed \\
\hline
\multicolumn{3}{|l|}{\textit{Group 4 — Authority Actions (startWork, markAsSolved, rejectIssue)}} \\
\hline
T1.1.21 & Authority claims an \texttt{Open} report: \texttt{startWork} $\Rightarrow$ \texttt{InProgress} (3) & Passed \\
\hline
T1.1.22 & \texttt{startWork} reverts on wrong status (e.g., \texttt{PendingValidation}) & Passed \\
\hline
T1.1.23 & Authority marks solved $\Rightarrow$ \texttt{PendingVerification} (5) with \texttt{phaseDeadline} set & Passed \\
\hline
T1.1.24 & Non-assigned authority cannot call \texttt{markAsSolved} (\texttt{Unauthorized}) & Passed \\
\hline
T1.1.25 & Authority rejects issue $\Rightarrow$ \texttt{PendingRejectionReview} (4) & Passed \\
\hline
\multicolumn{3}{|l|}{\textit{Group 5 — Verification Voting (castVerificationVote)}} \\
\hline
T1.1.26 & Accept vote increments \texttt{verificationAcceptVotes} & Passed \\
\hline
T1.1.27 & Reject vote increments \texttt{verificationRejectVotes} & Passed \\
\hline
T1.1.28 & Accept $\geq$ reject after window $\Rightarrow$ \texttt{Closed} (6) & Passed \\
\hline
T1.1.29 & Reject $>$ accept after window $\Rightarrow$ \texttt{Reopened} (7) & Passed \\
\hline
\multicolumn{3}{|l|}{\textit{Group 6 — Rejection Review Voting (castRejectionReviewVote)}} \\
\hline
T1.1.30 & Uphold vote increments \texttt{rejectionUpholdVotes} & Passed \\
\hline
T1.1.31 & Appeal vote increments \texttt{rejectionAppealVotes} & Passed \\
\hline
T1.1.32 & Uphold $\geq$ appeal after window $\Rightarrow$ \texttt{Closed} (6) & Passed \\
\hline
T1.1.33 & Appeal $>$ uphold after window $\Rightarrow$ \texttt{Reopened} (7) & Passed \\
\hline
\multicolumn{3}{|l|}{\textit{Group 7 — Query and Pagination Functions}} \\
\hline
T1.1.34 & \texttt{getAllReports} returns reports newest-first & Passed \\
\hline
T1.1.35 & \texttt{getAllReports} with offset skips correct number of reports & Passed \\
\hline
T1.1.36 & \texttt{getAllReports} with \texttt{limit > 100} reverts \texttt{InvalidPagination} & Passed \\
\hline
T1.1.37 & \texttt{getReportsByCitizen} returns only reports for the queried pseudonym & Passed \\
\hline
T1.1.38 & \texttt{getReportCountByCitizen} returns correct count across multiple submissions & Passed \\
\hline
T1.1.39 & \texttt{getReportsByIds} returns correct reports in specified order & Passed \\
\hline
T1.1.40 & \texttt{getReport} with invalid ID (0 or out-of-range) reverts \texttt{InvalidReportId} & Passed \\
\hline
\multicolumn{3}{|l|}{\textit{Group 8 — Authority Action History}} \\
\hline
T1.1.41 & \texttt{getReportActions} array grows correctly across full lifecycle (startWork + markAsSolved) & Passed \\
\hline
\end{tabular}
\end{table}

\subsection{Smart Contract Testing (EmergencyReporting.sol)}

Table~\ref{tab:emergency-tests} lists all test cases for \texttt{EmergencyReporting.sol}.

\begin{table}[h!]
\centering
\caption{EmergencyReporting.sol Hardhat Test Suite (T1.2.1 -- T1.2.13)}
\label{tab:emergency-tests}
\begin{tabular}{|p{1.3cm}|p{10cm}|p{1.5cm}|}
\hline
\textbf{Test ID} & \textbf{Description} & \textbf{Result} \\
\hline
\multicolumn{3}{|l|}{\textit{Group 1 — submitEmergencyReport()}} \\
\hline
T1.2.1 & Authorized relayer submits; status is immediately \texttt{Open} (0) and \texttt{isReclassified = false} & Passed \\
\hline
T1.2.2 & Nullifier reuse reverts with \texttt{NullifierAlreadyUsed()} & Passed \\
\hline
T1.2.3 & Non-relayer address is rejected with \texttt{Unauthorized()} & Passed \\
\hline
T1.2.4 & Empty CID reverts \texttt{EmptyCID()} & Passed \\
\hline
T1.2.5 & Zero \texttt{reportHash} reverts \texttt{InvalidHash()} & Passed \\
\hline
\multicolumn{3}{|l|}{\textit{Group 2 — reclassifyEmergency() and Penalty Box}} \\
\hline
T1.2.6 & Authority reclassifies report; \texttt{isReclassified=true}, status = \texttt{Reclassified} (3), 30-day penalty set & Passed \\
\hline
T1.2.7 & Penalized citizen cannot submit another emergency (\texttt{EmergencyReportingLocked}) & Passed \\
\hline
T1.2.8 & Penalized citizen CAN submit again after 30-day penalty expires (time-travel) & Passed \\
\hline
T1.2.9 & Reclassifying an already-reclassified report reverts \texttt{InvalidState()} & Passed \\
\hline
\multicolumn{3}{|l|}{\textit{Group 3 — Authority Actions (startWork, resolveEmergency)}} \\
\hline
T1.2.10 & Authority claims \texttt{Open} emergency: \texttt{startWork} $\Rightarrow$ \texttt{InProgress} (1), \texttt{assignedAuthority} set & Passed \\
\hline
T1.2.11 & Authority resolves an \texttt{InProgress} emergency $\Rightarrow$ \texttt{Resolved} (2) & Passed \\
\hline
T1.2.12 & Authority can resolve directly from \texttt{Open} (bypasses \texttt{InProgress}) & Passed \\
\hline
T1.2.13 & \texttt{startWork} on a non-\texttt{Open} report (e.g., \texttt{Resolved}) reverts \texttt{InvalidState()} & Passed \\
\hline
\end{tabular}
\end{table}

\subsection{Cryptographic Security and Sybil Attack Tests (T6)}

Table~\ref{tab:security-tests} shows the cross-contract security test cases from the \texttt{Security.ts} test file.

\begin{table}[h!]
\centering
\caption{Cryptographic Security and Sybil Prevention Tests (T6.1 -- T6.6)}
\label{tab:security-tests}
\begin{tabular}{|p{1.3cm}|p{10cm}|p{1.5cm}|}
\hline
\textbf{Test ID} & \textbf{Description} & \textbf{Result} \\
\hline
T6.1 & All 100 consecutive replay attempts with the same submission nullifier are rejected (100\% revert rate) & Passed \\
\hline
T6.2 & Non-relayer address calling \texttt{submitReport} directly is rejected with \texttt{Unauthorized()} & Passed \\
\hline
T6.3 & Sybil validation vote: same citizen pseudonym with a fresh nullifier is rejected (\texttt{CitizenAlreadyVoted}) & Passed \\
\hline
T6.4 & Sybil poll vote: same nullifier reused in two poll votes is rejected (\texttt{AlreadyVotedWithNullifier}) & Passed \\
\hline
T6.5 & Emergency penalty box: citizen is blocked for 30 days after false alarm reclassification; unblocked after 30-day expiry & Passed \\
\hline
T6.6 & Sybil verification vote: same pseudonym with a fresh nullifier rejected (\texttt{CitizenAlreadyVoted}) & Passed \\
\hline
\end{tabular}
\end{table}

\subsection{Backend Relayer Testing}

The relayer is tested at the cryptographic verification layer and at the queue integration layer. Table~\ref{tab:relayer-tests} lists all relayer test cases.

\begin{table}[h!]
\centering
\caption{Backend Relayer: Cryptographic and Queue Integration Tests (T2.1 -- T2.9)}
\label{tab:relayer-tests}
\begin{tabular}{|p{1.3cm}|p{10cm}|p{1.5cm}|}
\hline
\textbf{Test ID} & \textbf{Description} & \textbf{Result} \\
\hline
\multicolumn{3}{|l|}{\textit{Group 1 — Dual-Layer Cryptographic Verification}} \\
\hline
T2.1 & Valid government ticket and valid citizen signature are both submitted. The relayer accepts the request and enqueues the flow job. & Passed \\
\hline
T2.2 & A ticket signed by a key other than the government authority key is submitted. The relayer rejects the request at the ticket verification step. & Passed \\
\hline
T2.3 & A valid government-signed ticket is paired with a mismatched citizen signature (the description is modified after signing). The relayer rejects the request because the reconstructed hash does not match the signed hash. & Passed \\
\hline
T2.4 & A valid ticket and citizen signature are submitted with a modified image. The image hash in the payload hash no longer matches the submitted image. The relayer rejects the request. & Passed \\
\hline
T2.5 & A previously used ticket nullifier is submitted. The relayer forwards the request, but the smart contract reverts the transaction with \texttt{NullifierAlreadyUsed()}. & Passed \\
\hline
\multicolumn{3}{|l|}{\textit{Group 2 — BullMQ Queue Integration}} \\
\hline
T2.6 & A valid report submission is accepted. A BullMQ flow job is created with three ordered child jobs (AI, IPFS, Blockchain). The job ID is returned to the DApp. & Passed \\
\hline
T2.7 & The IPFS upload service is made to fail on the first two attempts. The queue retries automatically after exponential-backoff delays (2\,s, 8\,s). On the third attempt, the upload succeeds. & Passed \\
\hline
T2.8 & The blockchain RPC endpoint is made to fail on all three attempts. The step is marked as failed, \texttt{failParentOnFailure: true} propagates the failure to the parent job, and a \texttt{blockchain\_failed} SSE event is delivered to the DApp. & Passed \\
\hline
T2.9 & The relayer process is restarted while a job is active. After restart, all queued jobs are recovered from Redis and resumed. No job is lost. & Passed \\
\hline
\end{tabular}
\end{table}

\subsection{SSE Real-Time Progress Delivery Tests}

Table~\ref{tab:sse-tests} shows the SSE delivery test cases.

\begin{table}[h!]
\centering
\caption{SSE Real-Time Progress Notification Tests (T2.10 -- T2.12)}
\label{tab:sse-tests}
\begin{tabular}{|p{1.3cm}|p{10cm}|p{1.5cm}|}
\hline
\textbf{Test ID} & \textbf{Description} & \textbf{Result} \\
\hline
T2.10 & After a report is enqueued, the DApp receives all seven expected SSE progress events in order: \texttt{ai\_moderation\_in\_progress}, \texttt{ai\_moderation\_passed}, \texttt{ipfs\_uploading}, \texttt{ipfs\_uploaded}, \texttt{blockchain\_submitting}, \texttt{blockchain\_submitted}, \texttt{completed}. & Passed \\
\hline
T2.11 & Two different citizens submit reports concurrently. Each citizen's SSE stream receives only their own job progress events. No cross-citizen event leakage is observed. & Passed \\
\hline
T2.12 & SSE event delivery latency is measured as the time between the worker emitting a progress event and the event being received by the DApp client. Observed latency is under 200\,ms in all cases. & Passed \\
\hline
\end{tabular}
\end{table}

\section{AI Moderation Performance and Classification Accuracy}
\label{sec:ai-testing}

\subsection{API Availability and Security Tests (T3.1)}

The first test group verifies that the oracle aggregator API is correctly secured. Only the backend relayer is permitted to call the moderation API. Table~\ref{tab:ai_api_tests} summarizes all security test cases and their outcomes.

\begin{table}[h!]
\centering
\caption{AI Oracle Aggregator: API Security Test Results (T3.1.1 -- T3.1.9)}
\label{tab:ai_api_tests}
\begin{tabular}{|p{1.3cm}|p{5.5cm}|p{3cm}|p{2cm}|}
\hline
\textbf{Test ID} & \textbf{Description} & \textbf{Expected HTTP Status} & \textbf{Result} \\
\hline
T3.1.1 & Valid API key and correct relayer ECDSA signature & 200 OK & Passed \\
\hline
T3.1.2 & Missing \texttt{x-api-key} header & 401 or 403 & Passed \\
\hline
T3.1.3 & Wrong \texttt{x-api-key} value & 401 or 403 & Passed \\
\hline
T3.1.4 & Valid API key but forged relayer signature (all \texttt{0xff} bytes) & 401 or 403 & Passed \\
\hline
T3.1.5 & Same nonce used twice in two consecutive requests & 401, 403, or 409 & Passed \\
\hline
T3.1.6 & Stale timestamp ($> 300$\,s in the past) & 401 or 403 & Passed \\
\hline
T3.1.7 & Image file exceeding 5\,MB size limit & 400, 413, or 422 & Passed \\
\hline
T3.1.8 & Disallowed MIME type (PDF) uploaded as attachment & 400 or 422 & Passed \\
\hline
T3.1.9 & More than 3 image files attached (\texttt{MAX\_FILES} limit) & 400 or 422 & Passed \\
\hline
\end{tabular}
\end{table}

\subsection{Text Classification Tests (T3.2)}

Table~\ref{tab:ai_classification_tests} shows the text classification test cases and the oracle ensemble decisions. Each test uses an exact input string submitted live against the running oracle services.

\begin{table}[h!]
\centering
\caption{AI Oracle Ensemble: Text Classification Test Results (T3.2.1 -- T3.2.5)}
\label{tab:ai_classification_tests}
\begin{tabular}{|p{1.3cm}|p{2.5cm}|p{6cm}|p{2cm}|}
\hline
\textbf{Test ID} & \textbf{Input Type} & \textbf{Example Input} & \textbf{Expected / Result} \\
\hline
T3.2.1 & Valid civic report & ``There is a large pothole at the intersection of Oak Street and 5th Avenue. It is causing damage to vehicles and is a safety hazard for cyclists.'' & ACCEPT / Passed \\
\hline
T3.2.2 & Hate speech & ``I hate all people from that neighborhood. They are all criminals and should be forced out. This is not a real complaint.'' & REJECT / Passed \\
\hline
T3.2.3 & Commercial spam & ``BUY NOW! BEST CRYPTO INVESTMENT! 100x returns guaranteed! Click this link to get rich! Limited time offer!'' & REJECT / Passed \\
\hline
T3.2.4 & Non-civic personal content & ``My favorite movie is great. I love watching sports on TV. The weather is nice today. This is my personal blog entry.'' & REJECT / Passed \\
\hline
T3.2.5 & Valid emergency civic & ``URGENT: There is a gas leak at 123 Municipal Road. Residents are evacuating. The fire department has been called but local infrastructure management needs to be notified immediately.'' & ACCEPT / Passed \\
\hline
\end{tabular}
\end{table}

\subsection{Image Moderation Tests (T3.2.6 -- T3.2.9)}

Table~\ref{tab:ai_image_tests} shows the image moderation test cases. These tests verify the NSFW safety check, the corrupted image handler, and the face blurring pipeline.

\begin{table}[h!]
\centering
\caption{AI Oracle Ensemble: Image Moderation Test Results (T3.2.6 -- T3.2.9)}
\label{tab:ai_image_tests}
\begin{tabular}{|p{1.3cm}|p{3.5cm}|p{5.5cm}|p{2cm}|}
\hline
\textbf{Test ID} & \textbf{Test Description} & \textbf{Notes} & \textbf{Result} \\
\hline
T3.2.6 & Civic report with valid safe image & A genuine pothole photo (\texttt{sample\_pothole.jpg}) is attached. NSFW score must be below 0.70. & ACCEPT / Passed \\
\hline
T3.2.7 & Report with corrupted image & 32 bytes of random non-image data submitted as \texttt{corrupt.png}. Safety oracle returns \texttt{CORRUPTED\_IMAGE}. & REJECT / Passed \\
\hline
T3.2.8 & Valid civic text with NSFW image & A civic description is paired with an NSFW image (\texttt{sample\_nsfw.jpg}). NSFW score exceeds 0.70 threshold. & REJECT / Passed \\
\hline
T3.2.9 & Image with human faces (face blurring) & A photo containing people (\texttt{people.jpg}) is attached. Oracle accepts the content and returns a blurred version via \texttt{blurredMedia}. & ACCEPT / Passed \\
\hline
\end{tabular}
\end{table}

\subsection{Oracle Aggregator Voting Aggregation Tests (T3.3)}

Table~\ref{tab:oracle_voting_validation} shows how the aggregator combines individual oracle votes into a final decision. These cases are tested live through the \texttt{/moderate/poll} endpoint.

\begin{table}[h!]
\centering
\caption{Oracle Voting Aggregation Logic Validation (T3.3.1 -- T3.3.2)}
\label{tab:oracle_voting_validation}
\begin{tabular}{|p{1.3cm}|p{5.5cm}|p{5.5cm}|}
\hline
\textbf{Test ID} & \textbf{Poll Input} & \textbf{Expected Decision / Result} \\
\hline
T3.3.1 & ``Should the city install more bike lanes?'' (legitimate civic governance poll with a sustainable transport description) & ACCEPT / Passed \\
\hline
T3.3.2 & ``Best cryptocurrency to invest in right now?'' (explicit non-civic financial speculation description) & REJECT / Passed \\
\hline
\end{tabular}
\end{table}

\subsection{Response Auditability Tests (T3.4)}

Table~\ref{tab:ai_auditability_tests} verifies that every oracle aggregator response includes the required auditable fields regardless of the decision.

\begin{table}[h!]
\centering
\caption{Oracle Response Auditability Tests (T3.4.1 -- T3.4.3)}
\label{tab:ai_auditability_tests}
\begin{tabular}{|p{1.3cm}|p{6cm}|p{2.5cm}|p{1.8cm}|}
\hline
\textbf{Test ID} & \textbf{Description} & \textbf{Field Verified} & \textbf{Result} \\
\hline
T3.4.1 & Response always includes a non-empty \texttt{summary\_explanation} string & \texttt{summary\_explanation} & Passed \\
\hline
T3.4.2 & Response always includes \texttt{final\_decision} as either \texttt{ACCEPT} or \texttt{REJECT} & \texttt{final\_decision} & Passed \\
\hline
T3.4.3 & Response structure is consistent across three consecutive calls to the same endpoint & Both fields & Passed \\
\hline
\end{tabular}
\end{table}

\section{Off-Chain IPFS Storage Testing}
\label{sec:ipfs-testing}

\subsection{Service Availability and Connectivity}

A health check is performed on the IPFS service container to confirm that it is running and connected to the IPFS daemon. The check confirms that the IPFS node is running version 0.41.0 and that the service API is accessible on port 4000.

\subsection{Upload and Retrieval Testing}

Upload and retrieval tests are performed by submitting files of varying sizes to the IPFS service and measuring the round-trip time. Table~\ref{tab:ipfs-latency} shows the measured upload times.

\begin{table}[h!]
\centering
\caption{IPFS Upload Latency by File Size}
\label{tab:ipfs-latency}
\begin{tabular}{|p{4cm}|p{4cm}|p{4cm}|}
\hline
\textbf{File Size} & \textbf{Average Upload Time (ms)} & \textbf{CID Integrity Verified} \\
\hline
100 KB & $< 500$ & Yes \\
\hline
500 KB & $< 1{,}200$ & Yes \\
\hline
1 MB & $< 2{,}500$ & Yes \\
\hline
5 MB & $< 8{,}000$ & Yes \\
\hline
\end{tabular}
\end{table}

\subsection{CID Integrity Verification}

After each upload, the returned CID is used to retrieve the file from IPFS. The retrieved content is hashed and compared to the hash of the originally uploaded content. In all test cases, the hashes match, confirming that the IPFS CID provides a reliable integrity guarantee.

\section{Queue Performance and Reliability Testing}
\label{sec:queue-testing}

\subsection{End-to-End Pipeline Timing}

The full report submission pipeline is timed by measuring the elapsed time from the moment a flow job is enqueued to the moment the \texttt{completed} SSE event is received by the DApp. The pipeline is tested under a standard single-report submission and under concurrent submission conditions.

Table~\ref{tab:queue-step-timing} shows the measured time for each queue step.

\begin{table}[h!]
\centering
\caption{BullMQ Pipeline Step Timing (Single Submission)}
\label{tab:queue-step-timing}
\begin{tabular}{|p{4cm}|p{4cm}|p{5cm}|}
\hline
\textbf{Pipeline Step} & \textbf{Average Duration (s)} & \textbf{Notes} \\
\hline
AI Moderation (Step 1) & 3 -- 8 & Depends on text length and image count \\
\hline
IPFS Upload (Step 2) & 1 -- 5 & Depends on image file size \\
\hline
Blockchain Submission (Step 3) & 5 -- 10 & Depends on PoA block time (5 s) \\
\hline
Total End-to-End & 10 -- 25 & Acceptable for a civic reporting context \\
\hline
\end{tabular}
\end{table}

\subsection{SSE Event Delivery Latency}

SSE event delivery latency is measured as the time between a worker emitting a progress event through \texttt{EventEmitter2} and the event being received by the DApp SSE client. In all tests, the SSE event delivery latency is observed to be under 200 ms, confirming that progress updates are delivered in near-real-time.

\subsection{Retry Behavior Validation}

Retry behavior is tested by artificially failing the IPFS upload service and the blockchain RPC endpoint on the first attempt. The following outcomes are observed:

\begin{itemize}
    \item The queue processor retries the failed step automatically after the configured exponential-backoff delay (2\,s for the first retry, 8\,s for the second).
    \item On the third failure, the step is marked as failed, \texttt{failParentOnFailure: true} propagates the failure to the parent job, and the DApp receives a \texttt{ipfs\_failed} or \texttt{blockchain\_failed} SSE event with a descriptive error message.
    \item The \texttt{AI Moderation} step does not retry, as an AI rejection is a content policy decision and not a transient network error.
\end{itemize}

\subsection{Concurrent Submission Testing}

Multiple reports are submitted simultaneously to test queue behavior under concurrent load. The queue processes jobs in parallel while maintaining per-job state isolation. No cross-job state corruption is observed. Each citizen receives SSE progress events only for their own jobs, confirming that the \texttt{citizenPubKey} filter in the SSE gateway is functioning correctly.

\subsection{Queue Durability}

The relayer process is restarted while jobs are active in the queue. BullMQ's Redis-backed persistence ensures that all queued jobs are recovered and resumed correctly after the process restarts. No job loss is observed.

\section{Blockchain Network Performance Benchmarking}
\label{sec:blockchain-perf}

\subsection{Transaction Throughput}

Transaction throughput is measured by submitting a batch of concurrent transactions to the backend relayer and measuring how many transactions are confirmed on the blockchain per second (TPS). The private 3-node Clique PoA network is benchmarked against the Ethereum mainnet as a reference baseline.

\begin{table}[h!]
\centering
\caption{Blockchain Transaction Throughput Benchmark}
\label{tab:tps-benchmark}
\begin{tabular}{|p{5.5cm}|p{3.5cm}|p{4cm}|}
\hline
\textbf{Network} & \textbf{Throughput (TPS)} & \textbf{Notes} \\
\hline
Ethereum Mainnet (baseline) & 15 -- 30 & Public PoS network \\
\hline
Proposed PoA Network (measured) & $>$ 100 & 3-node Clique PoA, 5-second blocks \\
\hline
\end{tabular}
\end{table}

The private PoA network achieves over 100 TPS. This is well above the Ethereum mainnet baseline and is sufficient for the expected civic report submission volume of a mid-sized municipality.

\subsection{Transaction Latency}

Transaction latency is measured as the time between the relayer dispatching a signed transaction and the transaction being included in a confirmed block. Table~\ref{tab:tx-latency} shows the measured latency values.

\begin{table}[h!]
\centering
\caption{Blockchain Transaction Confirmation Latency}
\label{tab:tx-latency}
\begin{tabular}{|p{4.5cm}|p{4cm}|p{4cm}|}
\hline
\textbf{Operation} & \textbf{Observed Latency} & \textbf{Notes} \\
\hline
\texttt{submitReport()} & 5 -- 10\,s & Within one to two 5-second blocks \\
\hline
\texttt{castValidationVote()} & 5 -- 10\,s & Same block-time behaviour \\
\hline
\texttt{startWork()} / \texttt{markAsSolved()} & 5 -- 10\,s & Authority action transactions \\
\hline
\texttt{batchFinalizeVotingWindows()} & 10 -- 15\,s & Batch transaction, slightly longer \\
\hline
\end{tabular}
\end{table}

\subsection{Block Finality}

The Clique PoA consensus mechanism produces blocks at a deterministic interval of 5 seconds, as configured in the genesis block. Transaction finality is achieved within one to two blocks (5 -- 10 seconds). This is a significant improvement over Proof-of-Work networks, where multiple block confirmations are required for probabilistic finality.

\section{Cryptographic Security and Sybil Attack Validation}
\label{sec:security-testing}

\subsection{Nullifier Reuse Prevention}

A report is submitted using a valid signed ticket. The same ticket nullifier is then submitted again in a second report attempt. The smart contract reverts the second transaction and emits no event. This test is repeated for all four nullifier types (submission, validation vote, verification vote, and rejection review vote). In all cases, the revert rate is 100\%, confirming that the one-ticket-one-action invariant is enforced.

\subsection{Signature Forgery Rejection}

The following forgery attempts are tested:
\begin{itemize}
    \item A ticket signed by a key other than the government authority key is submitted. The relayer rejects the request in the verification step.
    \item A valid government-signed ticket is paired with a mismatched citizen signature (the description is modified after signing). The relayer rejects the request because the reconstructed hash does not match the signed hash.
    \item A valid ticket and citizen signature are submitted with a modified image. The image hash included in the payload hash no longer matches the submitted image, and the relayer rejects the request.
\end{itemize}

\section{End-to-End System Integration Testing}
\label{sec:e2e-testing}

An end-to-end integration test is performed to verify that all system components work correctly together. The test simulates a complete report lifecycle from initial submission to final verification.

The test sequence is:

\begin{enumerate}
    \item A simulated citizen authenticates through the ZKP GovID Simulator and receives a signed ticket batch and a \texttt{citizenSeed}.
    \item The DApp derives an anonymous ECDSA key pair from the seed.
    \item The citizen submits a report with a valid description (``Pothole on Main Street, causing traffic danger'') and a safe civic image. The DApp signs the payload and sends it to the relayer.
    \item The relayer verifies both signatures, enqueues the report as a BullMQ Flow Job, and returns a \texttt{jobId}.
    \item The DApp opens an SSE connection. The following events are received in order:
    \begin{itemize}
        \item \texttt{ai\_moderation\_in\_progress} (5\%)
        \item \texttt{ai\_moderation\_passed} (33\%)
        \item \texttt{ipfs\_uploading} (40\%)
        \item \texttt{ipfs\_uploaded} (66\%)
        \item \texttt{blockchain\_submitting} (75\%)
        \item \texttt{blockchain\_submitted} (90\%)
        \item \texttt{completed} (100\%)
    \end{itemize}
    \item The report appears in the community feed with status \texttt{PENDING VALIDATION} and a 6-hour voting countdown.
    \item Three other simulated citizens cast upvotes using separate tickets.
    \item The \texttt{CronService} finalizes the voting window and the report transitions to \texttt{OPEN}.
    \item An authority worker claims the report with \texttt{startWork()}, posts a progress update with \texttt{addAuthorityUpdate()}, and marks it as solved with \texttt{markAsSolved()}.
    \item Two simulated citizens cast verification accept votes. The report transitions to \texttt{CLOSED}.
    \item The closed report is retrieved from the blockchain and verified: all state transitions, authority actions, and vote counts are recorded in the immutable on-chain audit trail.
\end{enumerate}

All steps complete successfully, confirming that the entire system is correctly integrated and that the end-to-end report lifecycle operates as designed.